% GENERATED FILE. Edit canonical lessons, then run npm run generate:books.
% canonical-source-hash: fnv1a64:462906939d6ce29f
% canonical-lessons: HI-C535-daur, HI-C535-phurti, HI-C535-phadna, HI-C535-medal, HI-C535-daurana

\chapter{Training Hard}
\label{ch:hi-training-hard}
\hlchaptermodality{535}

\begin{chapteropening}
\textbf{By the end of this chapter:} \emph{I can talk about rounds, speed and scores.}

Complete the last lesson of chapter 535: I can talk about rounds, speed and scores.
\end{chapteropening}

\section[\texorpdfstring{\hi{दौर} (daur)}{दौर (daur)}]{\hi{दौर} (daur) — a round, a phase}

\label{lesson:HI-C535-daur}



\begin{quote}
Before the new one: say the Hindi for a winner, then the Hindi for a prize, an award.
\end{quote}

\subsection*{You'll want to know: \hi{दौर}}
\textbf{\hi{दौर}} — \emph{daur} — \textquotedblleft{}a round, a phase\textquotedblright{}.

\begin{quote}
\hi{पहले} \hi{दौर} \hi{में} \hi{हमारी} \hi{टीम} \hi{जीती।}
\end{quote}

\begin{quote}
\emph{Pahle daur meṁ hamārī ṭīm jītī.}
\end{quote}

\begin{quote}
Our team won in the first round.
\end{quote}

Masculine.

\begin{cousinweb}
\textbf{\hi{दौर}} (\emph{daur}) is usually traced to Arabic \emph{dawr}, \textquotedblleft{}a turn, a round\textquotedblright{}, through Persian.
\end{cousinweb}

\begin{grammarlens}[title={What you've built}]
One more word for describing what you lived and felt.
\end{grammarlens}

\subsection*{Your turn}
\begin{itemize}
\raggedright
  \item \emph{Say it:} \emph{daur}
  \item \emph{Say it:} \emph{daur}, once more
  \item \emph{Say it:} \emph{Pahle daur meṁ hamārī ṭīm jītī.} — Our team won in the first round.
  \item \emph{Recall:} say the Hindi for a winner, then the Hindi for a prize, an award, then say \emph{daur} again
\end{itemize}

\begin{tcolorbox}[breakable,colback=teal!4,colframe=teal!35!black,title={Before you move on}]
What does \hi{दौर} mean? (\textquotedblleft{}A round, a phase\textquotedblright{}.) Say it once more.
\end{tcolorbox}
\section[\texorpdfstring{\hi{फुर्ती} (phurtī)}{फुर्ती (phurtī)}]{\hi{फुर्ती} (phurtī) — quickness; agility}

\label{lesson:HI-C535-phurti}



\begin{quote}
Before the new one: say the Hindi for a prize, an award, then the Hindi for a round, a phase.
\end{quote}

\subsection*{You'll want to know: \hi{फुर्ती}}
\textbf{\hi{फुर्ती}} — \emph{phurtī} — \textquotedblleft{}quickness; agility\textquotedblright{}.

\begin{quote}
\hi{खिलाड़ी} \hi{ने} \hi{बड़ी} \hi{फुर्ती} \hi{से} \hi{गेंद} \hi{पकड़ी।}
\end{quote}

\begin{quote}
\emph{Khilāṛī ne baṛī phurtī se gẽd pakṛī.}
\end{quote}

\begin{quote}
The player caught the ball with great quickness.
\end{quote}

Feminine.

\begin{cousinweb}
\textbf{\hi{फुर्ती}} (\emph{phurtī}) is a tadbhava, inherited from Sanskrit \emph{sphurati}, \textquotedblleft{}darts, flashes\textquotedblright{} (CDIAL 13849). CDIAL gives the Hindi word as \textquotedblleft{}activity\textquotedblright{}; quickness is how Hindi uses it now.
\end{cousinweb}

\begin{grammarlens}[title={What you've built}]
One more word for describing what you lived and felt.
\end{grammarlens}

\subsection*{Your turn}
\begin{itemize}
\raggedright
  \item \emph{Say it:} \emph{phurtī}
  \item \emph{Say it:} \emph{phurtī}, once more
  \item \emph{Say it:} \emph{Khilāṛī ne baṛī phurtī se gẽd pakṛī.} — The player caught the ball with great quickness.
  \item \emph{Recall:} say the Hindi for a prize, an award, then the Hindi for a round, a phase, then say \emph{phurtī} again
\end{itemize}

\begin{tcolorbox}[breakable,colback=teal!4,colframe=teal!35!black,title={Before you move on}]
What does \hi{फुर्ती} mean? (\textquotedblleft{}Quickness; agility\textquotedblright{}.) Say it once more.
\end{tcolorbox}
\section[\texorpdfstring{\hi{फाँदना} (phā̃dnā)}{फाँदना (phā̃dnā)}]{\hi{फाँदना} (phā̃dnā) — to leap over; to vault}

\label{lesson:HI-C535-phadna}



\begin{quote}
Before the new one: say the Hindi for a round, a phase, then the Hindi for quickness.
\end{quote}

\subsection*{You'll want to know: \hi{फाँदना}}
\textbf{\hi{फाँदना}} — \emph{phā̃dnā} — \textquotedblleft{}to leap over; to vault\textquotedblright{}.

\begin{quote}
\hi{घोड़ा} \hi{ऊँची} \hi{दीवार} \hi{फाँद} \hi{गया।}
\end{quote}

\begin{quote}
\emph{Ghoṛā ū̃cī dīvār phā̃d gayā.}
\end{quote}

\begin{quote}
The horse leapt over the high wall.
\end{quote}

Past (he …): \hi{वह} \hi{नाली} \hi{फाँद} \hi{गया}, \emph{vah nālī phā̃d gayā}.

\begin{cousinweb}
\textbf{\hi{फाँदना}} (\emph{phā̃dnā}) is a tadbhava, inherited from Sanskrit \emph{spandatē}, \textquotedblleft{}quivers, throbs\textquotedblright{} (CDIAL 13806). CDIAL gives the Hindi verb as \textquotedblleft{}to leap\textquotedblright{}.
\end{cousinweb}

\begin{grammarlens}[title={What you've built}]
One more word for describing what you lived and felt.
\end{grammarlens}

\subsection*{Your turn}
\begin{itemize}
\raggedright
  \item \emph{Say it:} \emph{phā̃dnā}
  \item \emph{Say it:} \emph{phā̃dnā}, once more
  \item \emph{Say it:} \emph{Ghoṛā ū̃cī dīvār phā̃d gayā.} — The horse leapt over the high wall.
  \item \emph{Recall:} say the Hindi for a round, a phase, then the Hindi for quickness, then say \emph{phā̃dnā} again
\end{itemize}

\begin{tcolorbox}[breakable,colback=teal!4,colframe=teal!35!black,title={Before you move on}]
What does \hi{फाँदना} mean? (\textquotedblleft{}To leap over; to vault\textquotedblright{}.) Say it once more.
\end{tcolorbox}
\section[\texorpdfstring{\hi{मेडल} (meḍal)}{मेडल (meḍal)}]{\hi{मेडल} (meḍal) — a medal}

\label{lesson:HI-C535-medal}



\begin{quote}
Before the new one: say the Hindi for quickness, then the Hindi for to leap over.
\end{quote}

\subsection*{You'll want to know: \hi{मेडल}}
\textbf{\hi{मेडल}} — \emph{meḍal} — \textquotedblleft{}a medal\textquotedblright{}.

\begin{quote}
\hi{दौड़} \hi{में} \hi{मीरा} \hi{को} \hi{सोने} \hi{का} \hi{मेडल} \hi{मिला।}
\end{quote}

\begin{quote}
\emph{Dauṛ meṁ Mīrā ko sone kā meḍal milā.}
\end{quote}

\begin{quote}
Mira won a gold medal in the race.
\end{quote}

Masculine.

\begin{cousinweb}
\textbf{\hi{मेडल}} (\emph{meḍal}) is a loan from English \emph{medal}.
\end{cousinweb}

\begin{grammarlens}[title={What you've built}]
One more word for describing what you lived and felt.
\end{grammarlens}

\subsection*{Your turn}
\begin{itemize}
\raggedright
  \item \emph{Say it:} \emph{meḍal}
  \item \emph{Say it:} \emph{meḍal}, once more
  \item \emph{Say it:} \emph{Dauṛ meṁ Mīrā ko sone kā meḍal milā.} — Mira won a gold medal in the race.
  \item \emph{Recall:} say the Hindi for quickness, then the Hindi for to leap over, then say \emph{meḍal} again
\end{itemize}

\begin{tcolorbox}[breakable,colback=teal!4,colframe=teal!35!black,title={Before you move on}]
What does \hi{मेडल} mean? (\textquotedblleft{}A medal\textquotedblright{}.) Say it once more.
\end{tcolorbox}
\section[\texorpdfstring{\hi{दौड़ाना} (dauṛānā)}{दौड़ाना (dauṛānā)}]{\hi{दौड़ाना} (dauṛānā) — to make run; to send running}

\label{lesson:HI-C535-daurana}



\begin{quote}
Before the new one: say the Hindi for to leap over, then the Hindi for a medal.
\end{quote}

\subsection*{You'll want to know: \hi{दौड़ाना}}
\textbf{\hi{दौड़ाना}} — \emph{dauṛānā} — \textquotedblleft{}to make run; to send running\textquotedblright{}.

\begin{quote}
\hi{हर} \hi{सुबह} \hi{हमारे} \hi{शिक्षक} \hi{हमें} \hi{मैदान} \hi{में} \hi{दौड़ाते} \hi{हैं।}
\end{quote}

\begin{quote}
\emph{Har subah hamāre śikṣak hamẽ maidān meṁ dauṛāte haiṁ.}
\end{quote}

\begin{quote}
Every morning our teacher makes us run on the field.
\end{quote}

Past, with \hi{ने} (he or she …): \hi{उसने} \hi{घोड़ा} \hi{दौड़ाया}, \emph{usne ghoṛā dauṛāyā}.

\begin{cousinweb}
\textbf{\hi{दौड़ाना}} (\emph{dauṛānā}) is the causative of \emph{dauṛnā}, \textquotedblleft{}to run\textquotedblright{}, which comes from Sanskrit \emph{dravati}, \textquotedblleft{}runs\textquotedblright{} (CDIAL 6624).
\end{cousinweb}

\begin{grammarlens}[title={What you've built}]
One more word for describing what you lived and felt.
\end{grammarlens}

\subsection*{Your turn}
\begin{itemize}
\raggedright
  \item \emph{Say it:} \emph{dauṛānā}
  \item \emph{Say it:} \emph{dauṛānā}, once more
  \item \emph{Say it:} \emph{Har subah hamāre śikṣak hamẽ maidān meṁ dauṛāte haiṁ.} — Every morning our teacher makes us run on the field.
  \item \emph{Recall:} say the Hindi for to leap over, then the Hindi for a medal, then say \emph{dauṛānā} again
\end{itemize}

\begin{tcolorbox}[breakable,colback=teal!4,colframe=teal!35!black,title={Before you move on}]
What does \hi{दौड़ाना} mean? (\textquotedblleft{}To make run; to send running\textquotedblright{}.) Say it once more.
\end{tcolorbox}
